\section{Related Work}
\label{sec:related}

\para{Localized knowledge editing.}
ROME models mid-layer feed-forward modules as associative memories and writes a fact through a rank-one update \cite{meng2022rome}; MEMIT distributes related updates across layers to edit many facts \cite{meng2023memit}. AlphaEdit instead projects updates into a null space estimated from preservation representations \cite{fang2025alphaedit}. KnowEdit and EasyEdit organize evaluation around reliability, generalization, locality, and portability and show that no editor dominates all dimensions \cite{zhang2024knowedit}. These methods primarily edit entity--relation facts. We study a computed fact and include an exact external intervention that gives a causal zero-weight-change baseline. We do not treat sparse or restricted updates as evidence of locality: causal localization can be weakly related to the best place to edit \cite{hase2023localization}, and sequential editing can damage neighboring facts and downstream capabilities \cite{gupta2024catastrophic}.

\para{Propagation, rehearsal, and belief depth.}
RippleEdits shows that direct factual success does not ensure logical or compositional propagation \cite{cohen2024ripple}. Qin et al. connect these uneven ripple effects to gradient similarity and document over-ripple, negation, and cross-lingual failures \cite{qin2025messy}. Standard masked fine-tuning with paraphrase augmentation and locality rehearsal is a strong practical editor \cite{gangadhar2024finetuning}; our intervention ladder isolates those two components under a matched budget. A separate line of work uses diverse synthetic documents to implant false facts \cite{wang2025sdf}. Belief-depth evaluations distinguish direct answers from remote generalization, robustness under challenge, spontaneous salience, and internal evidence \cite{slocum2025believe}. Building on these distinctions, we report semantic transfer and untrained logical consequences separately. This choice is central: answering held-out paraphrases with \texttt{5} need not mean that the model uses the counterfactual consistently in arithmetic.

\para{Mechanisms of arithmetic.}
Mechanistic studies find that a small set of attention heads transmits operands and operators while downstream multilayer perceptrons promote answer tokens \cite{zhang2024arithmetic}. Arithmetic behavior also depends on overlapping heuristic neurons associated with operand ranges, result ranges, modular structure, and repeated operands \cite{nikankin2025arithmetic}. Complementary work identifies linear and periodic components in numerical representations for addition \cite{kantamneni2025trigonometry}. These results predict structured collateral effects: an edit may move inputs that share operands, outputs, token patterns, or gradients. We test those predictions behaviorally and through gradient similarity, but we do not infer a complete arithmetic circuit from correlation.

\para{Position of this work.}
Prior work establishes that edits can fail to propagate or can damage other behavior. Our contribution is narrower and complementary: we hold one computed counterfactual fixed, compare an exact exception with parameter updates, and map the boundary between intended semantic transfer, desired logical consequences, and structured collateral change. This design makes the phrase ``nothing else changed'' testable within a declared evaluation scope.
